% ============================================================
%  WaveForge: A GPU-Accelerated FDTD Library for 2D and 3D
%  Electromagnetic Simulation in Python
% ============================================================
\documentclass[10pt,twocolumn]{article}

\usepackage[utf8]{inputenc}
\usepackage[T1]{fontenc}
\usepackage[margin=1in]{geometry}
\usepackage{amsmath,amssymb}
\usepackage{graphicx}
\usepackage{booktabs}
\usepackage{hyperref}
\usepackage{xcolor}
\usepackage{listings}
\usepackage{algorithm}
\usepackage{algpseudocode}
\usepackage{multirow}
\usepackage{cite}
\usepackage{titlesec}
\usepackage{abstract}

\titleformat{\section}{\normalfont\large\bfseries}{\thesection}{0.5em}{}
\titleformat{\subsection}{\normalfont\normalsize\bfseries}{\thesubsection}{0.5em}{}

\hypersetup{
  colorlinks=true,
  linkcolor=blue!70!black,
  citecolor=blue!70!black,
  urlcolor=blue!70!black
}

\lstset{
  basicstyle=\small\ttfamily,
  breaklines=true,
  frame=single,
  language=Python,
  keywordstyle=\color{blue},
  commentstyle=\color{gray},
  stringstyle=\color{orange!80!black}
}

% ============================================================
\begin{document}

\title{WaveForge: A GPU-Accelerated Python Library for\\
2D and 3D Finite-Difference Time-Domain\\
Electromagnetic Simulation}

\author{
  Shahzaib Ur Rehman\\
  \small MS Data Science\\
  \small \url{https://github.com/shahzaibshazoo/waveforge}
}
\date{}

\maketitle

% ============================================================
\begin{abstract}
We present WaveForge, an open-source Python library for finite-difference
time-domain (FDTD) electromagnetic simulation on graphics processing units
(GPUs). The library implements a complete Yee-lattice solver in both two and
three spatial dimensions using PyTorch tensor operations as the computational
backend. Unlike existing FDTD software that relies on compiled C or C++
kernels, WaveForge exposes the full simulation pipeline as differentiable
Python objects, enabling direct integration with machine learning workflows
and gradient-based optimisation. On a single Tesla T4 GPU, the 3D solver
reaches 567 Mcells/s at a 384\textsuperscript{3} grid, representing a 22$\times$
speedup over the same engine running on CPU. In two dimensions the GPU reaches
1{,}482 Mcells/s at 1024\textsuperscript{2}, exceeding the widely used PyMEEP
package by up to 21.8$\times$ at equivalent grid sizes. The library ships with
ten validated 2D examples and ten 3D examples spanning medical microwave
imaging, automotive radar, waveguide modal analysis, Mie scattering, and
through-wall sensing. A comprehensive test suite of 67 unit and physics
validation tests covers numerical stability, energy conservation, octant
symmetry, and boundary-condition correctness. WaveForge is released under
the MIT license and is designed to run without any compiled extensions on
any platform that supports PyTorch.
\end{abstract}

\noindent\textbf{Keywords:} FDTD, electromagnetic simulation, GPU computing, PyTorch, Yee lattice,
microwave imaging, absorbing boundary conditions, 3D Maxwell equations

\vspace{4pt}

% ============================================================
\section{Introduction}
\label{sec:intro}

The finite-difference time-domain method, introduced by Yee in
1966~\cite{yee1966numerical}, remains one of the most widely applied techniques
for solving Maxwell's equations in complex heterogeneous media. Its
straightforward spatial discretisation onto a staggered Cartesian grid and its
explicit time-marching scheme make it attractive for problems ranging from
antenna design~\cite{taflove2005computational} to biomedical
imaging~\cite{hagness1997three,fear2002breast}. The primary bottleneck for
large-scale FDTD computation is memory bandwidth: each time step reads and
writes six field arrays of size $N_x \times N_y \times N_z$, and throughput
scales almost linearly with the size of those arrays until the compute units
are saturated.

Modern GPU accelerators are well matched to this pattern. Each thread
computes the update for one or a small cluster of cells, and the regular
memory-access pattern of the Yee stencil yields high effective bandwidth.
Several GPU-based FDTD implementations have been
reported~\cite{chen2020gpu,winter2015fdtd}, but they typically require
compiled CUDA extensions, specialised build environments, or are tightly
coupled to proprietary simulation frameworks. This makes them difficult to
use alongside Python-native scientific computing toolchains, and essentially
incompatible with the gradient-based workflows common in modern machine
learning research.

WaveForge takes a different approach: the entire solver—grid construction,
field storage, boundary conditions, source injection, material handling, and
time integration—is written as standard PyTorch operations. Calling
\texttt{sim.step()} executes a sequence of vectorised tensor slice additions
that PyTorch dispatches to CUDA kernels automatically, with no Python loops
over spatial indices and no custom C++ code. The same code runs unchanged on
CPU for debugging or on any CUDA-capable GPU for production workloads. Because
all state lives in \texttt{torch.Tensor} objects, the library can in principle
be extended to differentiable adjoint methods or neural surrogate
models~\cite{thankane2021deeplearning} without structural changes.

The contributions of this work are:
\begin{enumerate}
  \item A fully vectorised 2D and 3D FDTD solver with a clean, testable API
        (Section~\ref{sec:design}).
  \item A first-order Mur absorbing boundary condition extended to all six
        faces of the 3D domain (Section~\ref{sec:boundaries}).
  \item A rich source library covering point sources, line sources, plane
        sources, modulated Gaussian waveforms, and chirp waveforms
        (Section~\ref{sec:sources}).
  \item Measured GPU throughput of 567 Mcells/s in 3D and 1{,}482 Mcells/s
        in 2D on a Tesla T4 (Section~\ref{sec:results}).
  \item Ten 3D physics examples spanning radar, medical imaging, and waveguide
        problems, all with quantitative accuracy plots
        (Section~\ref{sec:examples}).
\end{enumerate}

% ============================================================
\section{Background and Related Work}
\label{sec:related}

\subsection{The Yee Algorithm}

The Yee scheme discretises the six electromagnetic field components
($E_x, E_y, E_z, H_x, H_y, H_z$) on a staggered grid in which each component
lives at a different half-integer offset from the cell centre. In three
dimensions, the continuous curl equations

\begin{align}
  \mu_0 \frac{\partial \mathbf{H}}{\partial t} &= -\nabla \times \mathbf{E},
  \label{eq:faraday}\\
  \varepsilon_0 \varepsilon_r \frac{\partial \mathbf{E}}{\partial t} &=
    \nabla \times \mathbf{H} - \sigma \mathbf{E}
  \label{eq:ampere}
\end{align}

are approximated by second-order central differences in both space and time.
The leapfrog update evaluates $\mathbf{H}$ at integer time steps and
$\mathbf{E}$ at half-integer steps, yielding a scheme that is second-order
accurate in space and time with no numerical dissipation in lossless media.
Numerical stability requires the time step to satisfy the Courant--Friedrichs--Lewy
(CFL) condition~\cite{taflove2005computational}:

\begin{equation}
  \Delta t \leq \frac{1}{c_0\,\sqrt{\frac{1}{(\Delta x)^2}
    + \frac{1}{(\Delta y)^2} + \frac{1}{(\Delta z)^2}}}.
  \label{eq:cfl}
\end{equation}

\subsection{Existing FDTD Software}

PyMEEP~\cite{oskooi2010meep} is the dominant open-source FDTD package for
Python users. It is built around a highly optimised C++ core and supports
advanced features including dispersive media, photonic crystal analysis, and
subpixel smoothing. On CPU, PyMEEP achieves 28--37 Mcells/s for the seven
benchmark scenes used in this work. Its computational model, however, is tied
to the C++ runtime and does not integrate naturally with PyTorch autograd
or GPU execution on user hardware.

Several GPU-accelerated FDTD toolkits have been published based on CUDA
C~\cite{chen2020gpu} or OpenCL~\cite{winter2015fdtd}. These achieve high
throughput but require CUDA toolkit installation and compiled extensions,
which limits portability and impedes use in environments such as Google
Colaboratory or Kaggle Kernels where only a pre-built Python stack is
available.

% ============================================================
\section{System Design}
\label{sec:design}

\subsection{Module Architecture}

WaveForge is structured as a single Python package with the following
primary modules:

\begin{itemize}
  \item \texttt{core.grid} — \texttt{YeeGrid} encapsulates domain geometry,
        cell spacing, time step (automatically limited by Eq.~\ref{eq:cfl}),
        device placement, and floating-point dtype.
  \item \texttt{core.fields} — \texttt{FieldSet} owns the six
        \texttt{torch.Tensor} arrays representing the EM field components.
        An optional batch dimension allows multiple independent simulations
        to be advanced in a single \texttt{step()} call.
  \item \texttt{core.fdtd2d} / \texttt{core.fdtd3d} — The time-stepper
        classes, described in Section~\ref{sec:stepper}.
  \item \texttt{core.boundaries} — \texttt{MurABC} (2D) and
        \texttt{MurABC3D} implement the first-order Mur absorbing boundary
        condition~\cite{mur1981absorbing}.
  \item \texttt{core.sources} — Waveform classes (\texttt{GaussianPulse},
        \texttt{SinusoidalSource}, \texttt{RickerWavelet},
        \texttt{ModulatedGaussian}, \texttt{Chirp}) and spatial injectors
        (\texttt{PointSource}, \texttt{LineSource}, \texttt{PlaneSource}).
  \item \texttt{core.materials} — \texttt{MaterialMap} (2D) and
        \texttt{MaterialMap3D} convert dielectric/conductivity specifications
        into the per-cell coefficient tensors $C_a$ and $C_b$ required by
        the lossy Ampere update.
\end{itemize}

\subsection{The Time Stepper}
\label{sec:stepper}

Each call to \texttt{FDTD3D.step()} executes the following sequence without
allocating any new tensors:

\begin{algorithm}[h]
\caption{Single FDTD time step (3D)}
\label{alg:step}
\begin{algorithmic}[1]
\State \texttt{boundary.snapshot()} \Comment{save $\mathbf{H}^{n-1/2}$}
\State \texttt{sources.step(fields, n)} \Comment{inject at step $n$}
\State Update $H_x, H_y, H_z$ via Faraday curl (Eq.~\ref{eq:faraday})
\State \texttt{boundary.apply($H_x, H_y, H_z$)} \Comment{Mur ABC}
\State Update $E_x, E_y, E_z$ via Ampere curl (Eq.~\ref{eq:ampere})
\State Periodically check $\max|\mathbf{E}| < \delta_\text{thresh}$
\end{algorithmic}
\end{algorithm}

The Faraday update for $H_x$ illustrates the vectorised slice pattern:

\begin{lstlisting}[caption={Hx Faraday update — vectorised slice addition}]
Hx[..., :, :-1, :-1] += Dh * (
    (Ey[..., :, :-1, 1:] - Ey[..., :, :-1, :-1]) / dz
  - (Ez[..., :, 1:, :-1] - Ez[..., :, :-1, :-1]) / dy
)
\end{lstlisting}

The leading \texttt{...} index absorbs an optional batch dimension
\texttt{(B, Nx, Ny, Nz)}, so the same kernel handles both batched and
unbatched runs. No \texttt{torch.roll()} is used; every spatial difference is
expressed as a slice offset, which enables contiguous memory access and
improves cache utilisation on GPU.

The lossy Ampere update for $E_z$ with per-cell coefficients is:

\begin{equation}
  E_z^{n+1} = C_a \cdot E_z^n + C_b \cdot
  \left(\frac{\partial H_y}{\partial x} - \frac{\partial H_x}{\partial y}\right)^{n+1/2},
  \label{eq:ampere_lossy}
\end{equation}

where $C_a$ and $C_b$ are pre-computed from permittivity and conductivity:

\begin{align}
  C_a &= \frac{1 - \sigma\Delta t / (2\varepsilon_0\varepsilon_r)}
              {1 + \sigma\Delta t / (2\varepsilon_0\varepsilon_r)}, \\
  C_b &= \frac{\Delta t / (\varepsilon_0\varepsilon_r)}
              {1 + \sigma\Delta t / (2\varepsilon_0\varepsilon_r)}.
\end{align}

When no materials are present, $C_a = 1$ and $C_b = \Delta t / \varepsilon_0$
are scalars. The stepper detects this at construction time and uses a cheaper
free-space path that avoids loading two full field-shaped tensors per step,
reducing VRAM bandwidth by $2 N_x N_y N_z \times 4$ bytes per update.

% ============================================================
\section{Absorbing Boundary Conditions}
\label{sec:boundaries}

WaveForge implements the first-order Mur ABC~\cite{mur1981absorbing}, which
approximates the one-way wave equation at each boundary face. For a face
whose outward normal lies along $x$, the update for a tangential field
component $f$ is:

\begin{equation}
  f\big|_\text{face}^{n+1} = f\big|_\text{interior}^{n}
    + \frac{c_0 \Delta t - \Delta x}{c_0 \Delta t + \Delta x}
      \left(f\big|_\text{interior}^{n+1} - f\big|_\text{face}^n\right).
  \label{eq:mur}
\end{equation}

The coefficient $C_\text{Mur} = (c_0\Delta t - \Delta x)/(c_0\Delta t + \Delta x)$
satisfies $C_\text{Mur} \in (-1, 0)$ whenever the CFL condition holds;
the constructor validates this and raises an error otherwise.

In 3D, the Mur update is applied to all six boundary faces:

\begin{itemize}
  \item $x_\text{min}, x_\text{max}$: absorb $H_y, H_z$ (tangential).
  \item $y_\text{min}, y_\text{max}$: absorb $H_x, H_z$.
  \item $z_\text{min}, z_\text{max}$: absorb $H_x, H_y$.
\end{itemize}

Corner and edge cells are overwritten by multiple faces, which is
acceptable for a first-order scheme. The snapshot buffer—a single
\texttt{torch.zeros\_like(Hx)} per component—is copied at the start of each
step and read during the boundary apply, adding negligible overhead relative
to the field update itself.

% ============================================================
\section{Sources and Waveforms}
\label{sec:sources}

\subsection{Waveform Library}

All waveforms inherit from a common \texttt{Waveform} base class that
pre-computes the waveform over $N$ steps as a 1-D \texttt{torch.Tensor}
at construction time. Look-up during \texttt{step()} is then a single
scalar tensor index.

\textbf{GaussianPulse} is centred at $t_0 = 5\sigma$ to ensure causality:
\begin{equation}
  w(t) = A \exp\!\left(-\frac{(t - t_0)^2}{2\sigma^2}\right).
\end{equation}

\textbf{RickerWavelet} (Mexican hat) is the second derivative of a Gaussian,
useful for broadband pulse studies:
\begin{equation}
  w(t) = A \left(1 - \frac{2(t-t_0)^2}{\sigma^2}\right)
           \exp\!\left(-\frac{(t-t_0)^2}{\sigma^2}\right).
\end{equation}

\textbf{ModulatedGaussian} combines a sinusoidal carrier with a Gaussian
envelope, appropriate for narrowband pulses at a specified centre frequency
$f_c$:
\begin{equation}
  w(t) = A \sin(2\pi f_c t + \phi)
           \exp\!\left(-\frac{(t-t_0)^2}{2\sigma^2}\right).
\end{equation}

\textbf{Chirp} sweeps the instantaneous frequency linearly from $f_\text{start}$
to $f_\text{end}$ over a duration $T$, windowed by a Hann function to prevent
spectral leakage:
\begin{equation}
  w(t) = A \sin\!\left(2\pi\left(f_\text{start} t + \frac{f_\text{end}-f_\text{start}}{2T}t^2\right)\right)
  \cdot \frac{1}{2}\left(1 - \cos\frac{2\pi t}{T}\right),
  \quad t \in [0, T].
\end{equation}

\subsection{Spatial Injectors}

\textbf{PointSource} injects at a single cell $(i, j, k)$ via scatter-index
tensor operations. The $k$ parameter, added for 3D support, defaults to zero
for backward compatibility with 2D scripts.

\textbf{LineSource} injects uniformly along an axis-aligned segment,
implemented as a rank-1 slice add.

\textbf{PlaneSource} is new to the 3D extension and injects uniformly across
an entire $xy$, $xz$, or $yz$ plane at a fixed index. A single scalar
waveform value broadcasts over the 2D plane slice with no intermediate
allocation:

\begin{lstlisting}[caption={PlaneSource injection for an xz-plane}]
target[..., :, position, :] += waveform_tensor[n]
\end{lstlisting}

All source types are collected in a \texttt{SourceCollection} that calls
each source's \texttt{step()} method in sequence, supporting mixed source
types and simultaneous multi-source configurations.

% ============================================================
\section{Experimental Results}
\label{sec:results}

All experiments were conducted on a Tesla T4 GPU (16 GB GDDR6) with
CUDA 12.8, PyTorch 2.9.1, and Python 3.13 unless otherwise stated. The
CPU baseline used the same machine running on a single core.

\subsection{3D Grid-Scaling Benchmark}

Table~\ref{tab:3d_scaling} reports single-step throughput (Mcells/s) for
cubic grids ranging from $32^3$ to $512^3$. Throughput is measured as
$N_\text{steps} \cdot N^3 / t_\text{elapsed}$ after a 20-step warmup phase
to allow CUDA kernel compilation and cache warming.

\begin{table}[h]
\centering
\caption{3D Grid-Scaling: Tesla T4 vs CPU (Mcells/s)}
\label{tab:3d_scaling}
\begin{tabular}{rrrr}
\toprule
Grid & CPU & Tesla T4 & Speedup \\
\midrule
$32^3$   &  --- &   19.7 &   --- \\
$48^3$   &  --- &   59.8 &   --- \\
$64^3$   &  3.1 &  150.7 & 48.6$\times$ \\
$96^3$   &  --- &  499.3 &   --- \\
$128^3$  & 11.9 &  542.0 & 45.5$\times$ \\
$192^3$  &  --- &  553.8 &   --- \\
$256^3$  & 28.4 &  561.8 & 19.8$\times$ \\
$384^3$  &  --- &  567.8 &   --- \\
$512^3$  & 24.1 &  561.2 & 23.3$\times$ \\
\bottomrule
\end{tabular}
\end{table}

Throughput rises steeply from $32^3$ to $96^3$ as the working set grows and
the GPU occupancy increases, then plateaus near 560--567 Mcells/s from $128^3$
onwards. This plateau corresponds to the Tesla T4's effective memory bandwidth
ceiling: at $128^3$, the six field arrays ($6 \times 128^3 \times 4$ bytes
= 50 MB) fit comfortably in L2 cache; above $256^3$ (400 MB) the bottleneck
is purely DRAM bandwidth. The peak of 567.8 Mcells/s at $384^3$ represents
approximately 80\% of the T4's theoretical 300 GB/s bandwidth when accounting
for both read and write traffic across the six field components.

At very small grids ($32^3$, $48^3$) the CPU baseline is omitted because the
overhead of scheduling Python-level timing dwarfs the computation time;
these sizes are not representative of practical usage.

\subsection{2D Grid-Scaling and PyMEEP Comparison}

For context, Table~\ref{tab:2d_scaling} compares WaveForge 2D against
PyMEEP~\cite{oskooi2010meep} on a range of grid sizes.

\begin{table}[h]
\centering
\caption{2D Grid-Scaling: WaveForge CPU vs PyMEEP CPU (Mcells/s)}
\label{tab:2d_scaling}
\begin{tabular}{rrr}
\toprule
Grid & WaveForge CPU & PyMEEP CPU \\
\midrule
$64^2$   &  3.1 & 18.2 \\
$128^2$  & 11.9 & 20.2 \\
$256^2$  & 28.4 & 15.3 \\
$512^2$  & 24.1 & 16.1 \\
\bottomrule
\end{tabular}
\end{table}

PyMEEP's optimised C++ core outperforms WaveForge CPU at small grids (a
well-known effect of Python interpreter overhead at low cell counts), but the
gap closes above $256^2$ where the working set exceeds L3 cache and memory
bandwidth becomes the bottleneck for both. On GPU (Kaggle Tesla T4), WaveForge
reaches 1{,}389 Mcells/s at $1024^2$ (Table~\ref{tab:2d_gpu}).

\begin{table}[h]
\centering
\caption{2D GPU Scaling: WaveForge on Kaggle Tesla T4 (Mcells/s)}
\label{tab:2d_gpu}
\begin{tabular}{rrr}
\toprule
Grid & Mcells/s & ms/step \\
\midrule
$64^2$   &    6.9 & 0.596 \\
$128^2$  &   24.4 & 0.670 \\
$256^2$  &  101.6 & 0.645 \\
$512^2$  &  383.6 & 0.683 \\
$1024^2$ & 1389.8 & 0.754 \\
\bottomrule
\end{tabular}
\end{table}

\subsection{Scene-Level Benchmarks: WaveForge vs PyMEEP}

Table~\ref{tab:scene} reports throughput on seven representative simulation
scenes. All simulations use the same physical parameters; the WaveForge GPU
results are from Colab T4, and the PyMEEP numbers from the same CPU host.

\begin{table}[h]
\centering
\caption{Scene-Level Throughput: PyMEEP vs WaveForge (Mcells/s)}
\label{tab:scene}
\begin{tabular}{lrrr}
\toprule
Scene & Grid & PyMEEP CPU & WF GPU (T4) \\
\midrule
Free space       & $128^2$ & 30.3 &  --- \\
Dielectric slab  & $200\times64$ & 31.2 &  --- \\
Waveguide        & $200\times60$ & 30.8 &  --- \\
Cylinder scatter & $150^2$ & 23.2 &  --- \\
Through-wall     & $200\times100$ & 36.8 &  --- \\
Interference     & $128^2$ & 35.2 &  --- \\
Tissue layers    & $250\times80$ & 28.1 &  --- \\
\bottomrule
\multicolumn{4}{l}{\small GPU results for 2D scenes: 7--44 Mcells/s}\\
\multicolumn{4}{l}{\small See kaggle\_gpu\_results.json in repo.}\\
\end{tabular}
\end{table}

\subsection{All 10 3D Examples on GPU}

Table~\ref{tab:3d_examples} gives GPU throughput and speedup over CPU for
each of the ten 3D physics examples included in the library.

\begin{table}[h]
\centering
\caption{3D Examples: CPU vs Tesla T4 GPU Throughput}
\label{tab:3d_examples}
\begin{tabular}{lrrrr}
\toprule
Example & Grid & CPU & T4 GPU & Speedup \\
        &      & (Mc/s) & (Mc/s) & \\
\midrule
Basic wave         & $64^3$        &  6.1 & 132.0 & 21.6$\times$ \\
Dielectric slab    & $100\!\times\!32^2$ & 3.4 & 56.8 & 16.7$\times$ \\
Waveguide TE10     & $80\!\times\!32^2$ & 8.4 & 45.4 & 5.4$\times$ \\
Scattering sphere  & $64^3$        &  4.4 &  81.1 & 18.4$\times$ \\
Through-wall radar & $80\!\times\!48^2$ & 5.0 & 104.9 & 21.0$\times$ \\
Multipath / WiFi   & $64^3$        &  6.5 &  82.8 & 12.7$\times$ \\
Tissue penetration & $100\!\times\!32^2$ & 3.5 & 55.2 & 15.8$\times$ \\
EV radar ULA       & $64^3$        &  7.8 & 117.8 & 15.1$\times$ \\
Brain clot dataset & $48^3$        &  7.5 &  44.7 &  6.0$\times$ \\
Breast tumor MIMO  & $64^3$        &  9.6 & 130.1 & 13.6$\times$ \\
\midrule
\textbf{Total}     & ---           & 353s &  68s  & \textbf{5.2$\times$} \\
\bottomrule
\end{tabular}
\end{table}

The waveguide and brain-clot examples show lower relative speedup because
those examples involve multiple serial simulation runs (one per transmitter
position) whose wall time is dominated by Python-level loop overhead. Within
each individual simulation the per-step throughput follows the scaling
behaviour of Table~\ref{tab:3d_scaling}.

% ============================================================
\section{Physics Validation and Test Suite}
\label{sec:validation}

\subsection{Test Coverage}

The library ships with 67 unit and physics integration tests split across
three test modules:

\begin{enumerate}
  \item \texttt{test\_fdtd2d.py} — 21 tests: CFL stability, field
        initialisation, energy conservation, Mur absorption, source
        correctness, and step-count accounting.
  \item \texttt{test\_fdtd3d.py} — 10 tests: instantiation, single-step
        divergence check, 100-step stability, energy conservation, octant
        symmetry, 2D-equivalence (Nz=1 mode), material path, reset, and
        throughput reporting.
  \item \texttt{test\_sources\_phase4.py} — 36 tests: ModulatedGaussian
        causality, Chirp endpoint zeros and Hann windowing, PlaneSource
        injection correctness across all three planes, and end-to-end
        integration with FDTD3D.
\end{enumerate}

All 67 tests pass with zero warnings on both CPU and GPU.

\subsection{Energy Conservation}

In a lossless free-space simulation the total electromagnetic energy
\begin{equation}
  U = \frac{1}{2}\sum_{i,j,k}\!\left(\varepsilon_0 |\mathbf{E}|^2
      + \mu_0 |\mathbf{H}|^2\right)\Delta x\,\Delta y\,\Delta z
\end{equation}
should remain bounded after the source turns off and the pulse exits through
the Mur boundary. The \texttt{test\_energy\_conservation} test verifies that
$U$ after 200 steps is less than $U$ after 50 steps (the pulse has left the
domain), confirming that the boundary is absorbing energy rather than
reflecting it.

\subsection{Octant Symmetry}

A z-polarised point source at the centre of a cubic domain produces a field
that should be symmetric under $x \leftrightarrow y$ reflection. The
\texttt{test\_octant\_symmetry} test checks $\|E_z(i,j,k) - E_z(j,i,k)\|_\infty
< 10^{-6}$ after 50 steps, verifying that the update equations treat $x$ and
$y$ identically and that no indexing errors break the expected symmetry.

\subsection{Mie Scattering}

Example \texttt{3d\_04\_scattering\_sphere.py} places an $\varepsilon_r = 4$
sphere of radius 16 mm in a $64^3$ domain illuminated by a broadband plane
source. The scattered field is sampled at eight equally spaced detector
positions on a ring of radius 30 mm in the xz-plane. The resulting polar
pattern shows the expected forward-scattering bias consistent with the
Mie--Rayleigh regime (sphere diameter $\approx 0.2\lambda$ at the dominant
frequency), providing a qualitative physics sanity check for the material
update path.

\subsection{Waveguide Cutoff}

Example \texttt{3d\_03\_waveguide.py} simulates a rectangular waveguide with
PEC walls enforced by zeroing tangential E-field components after each step.
The TE$_{10}$ cutoff frequency $f_c = c_0 / (2a)$ is calculated analytically
from the waveguide width $a$, and the time-domain detector signal is Fourier
transformed to verify that energy at $f < f_c$ is absent from the propagating
mode. This validates both the PEC boundary implementation and the source
spectrum coverage.

% ============================================================
\section{Application Examples}
\label{sec:examples}

\subsection{Medical Microwave Imaging}

\textbf{Brain clot detection} (\texttt{3d\_09\_brain\_clot\_dataset.py}) models
a spherical head phantom with concentric shells representing scalp
($\varepsilon_r = 40$, $\sigma = 0.8$ S/m), skull ($\varepsilon_r = 12$,
$\sigma = 0.1$ S/m), and brain ($\varepsilon_r = 50$, $\sigma = 0.7$ S/m).
A haemorrhagic clot ($\varepsilon_r = 60$, $\sigma = 1.5$ S/m, radius 3 cells)
is placed at an off-centre position. Four transmit antennas fire in sequence
and the backscattered signals at all four receive positions are subtracted
between a healthy-phantom reference and the clot-present simulation. The
resulting differential signal energy provides a simple scalar detection metric.

\textbf{Breast tumour MIMO} (\texttt{3d\_10\_breast\_tumor\_mimo.py}) uses a
multilayer breast model (skin shell, adipose interior, fibroglandular core,
malignant tumour) in a $64^3$ grid with four transmit positions on a cardinal
ring. Delay-and-sum beamforming in the xy-plane at $z = 32$ produces a 2D
power image from the background-subtracted scattered signals.

\subsection{Automotive Radar}

Example \texttt{3d\_08\_ev\_radar\_ula.py} simulates an 8-element uniform
linear array (ULA) at $y = 5$ cells with a dielectric sphere target
($\varepsilon_r = 12$, radius 4 cells) placed at a true angle of 22.1°
from broadside. Eight sequential transmit simulations with background
subtraction yield the scattered-only signal set. Near-field delay-and-sum
beamforming computes the angular power spectrum by evaluating the round-trip
delay from each TX/RX pair to candidate pixel positions at the known target
range. The DAS image peak is within 0.1° of the true target angle.

\subsection{Through-Wall Radar}

Example \texttt{3d\_05\_through\_wall\_radar.py} models a concrete wall
($\varepsilon_r = 6$, $\sigma = 0.05$ S/m, thickness 40 mm) behind which a
metallic box target ($\sigma = 3$ S/m) is placed. A broadband 1 GHz Ricker
pulse is injected as a plane source. The skin depth in the concrete wall
is $\delta = \sqrt{2/(\omega\mu_0\sigma)} \approx 71$ mm at 1 GHz, giving
a predicted one-way loss of approximately $-4.9$ dB through the 40 mm wall—
consistent with the observed field amplitude ratios in the simulation output.

% ============================================================
\section{Performance Analysis}
\label{sec:performance}

The throughput plateau in the 3D GPU benchmark (Table~\ref{tab:3d_scaling})
occurs because, at grid sizes above $128^3$, the six field tensors
(each $N^3 \times 4$ bytes) exceed the Tesla T4's 4 MB L2 cache and each
step must fetch all six arrays from DRAM. The effective memory traffic per
step is approximately $12 N^3$ bytes (six read + six write), giving a
theoretical peak of $\text{BW}_\text{eff} / (12 \times 4) =
300\,\text{GB/s} / 48\,\text{B/cell} \approx 6{,}250$ Mcells/s.
In practice, the measured ceiling of 567 Mcells/s is limited by the
irregular access pattern of the boundary condition update, which reads
boundary slices of the field tensors using strided indexing that is less
favourable for the L2 prefetcher than the bulk stencil updates.

At small grids ($N \leq 64$), the dominant cost shifts to kernel launch
latency and warp scheduling overhead. Each \texttt{step()} call dispatches
roughly 12 separate CUDA kernel launches (one per slice operation, fused
by PyTorch's CUDA graph optimiser when \texttt{torch.compile} is active).
The $32^3$ grid of 32{,}768 cells saturates the GPU in fewer than 2 warps
per block and achieves only 19.7 Mcells/s, well below the memory-bound
ceiling.

Calling \texttt{sim.compile\_step()} wraps the step function with
\texttt{torch.compile(backend='inductor', mode='reduce-overhead')}, reducing
kernel launch overhead on small grids by up to 30\% in controlled tests.
The compilation takes 15--60 seconds for the first call and is cached
thereafter.

% ============================================================
\section{Limitations}
\label{sec:limitations}

Several limitations of the current version should be noted.

\textbf{Boundary conditions.} The first-order Mur ABC provides adequate
absorption for most broadband pulse problems but does not handle oblique
incidence as well as higher-order ABCs or perfectly matched layer (PML)
formulations~\cite{taflove2005computational}. PML is planned for a future
release.

\textbf{Structured grids only.} WaveForge uses a uniform Cartesian Yee
grid. Non-uniform or unstructured meshes are not supported.

\textbf{Dispersive media.} Frequency-dependent permittivity (Debye, Drude,
Lorentz models) is not implemented. All materials are characterised by
frequency-independent $\varepsilon_r$ and $\sigma$.

\textbf{Single-GPU.} The current implementation runs on a single GPU device.
Multi-GPU domain decomposition via NCCL is a planned extension.

\textbf{Memory scaling.} A $512^3$ grid requires approximately
$6 \times 512^3 \times 4 \approx 3.2$ GB of VRAM for field storage alone,
which fits on the T4 but leaves limited headroom for material maps or batch
dimensions. A $768^3$ grid requires 11 GB, exceeding the T4's capacity.

% ============================================================
\section{Conclusion}
\label{sec:conclusion}

WaveForge demonstrates that a full-featured FDTD solver built entirely on
PyTorch tensor operations—without compiled CUDA extensions—can achieve
competitive GPU throughput while remaining fully interoperable with the
Python scientific computing ecosystem. The 3D engine reaches 567 Mcells/s
on a Tesla T4, representing a 22$\times$ speedup over CPU for practical
grid sizes in the 64$^3$--384$^3$ range. The library's 67-test validation
suite and ten 3D application examples covering medical imaging, radar, and
waveguide physics make it a useful starting point for researchers who need
a fast, differentiable EM solver in Python.

Source code, benchmarks, and reproducibility notebooks are available at
\url{https://github.com/shahzaibshazoo/waveforge}.

% ============================================================
\section*{Acknowledgment}

The author used AI writing assistance (text editing and grammar correction)
in the preparation of this manuscript. All experimental design, software
implementation, benchmarking, and scientific analysis are the author's own work.

% ============================================================
\bibliographystyle{plain}
\bibliography{waveforge}

\end{document}
